% !PRINTED 189
\weberpage{189}

Daraus läßt sich $a_{m,n}$ finden, wenn die früheren $a_{m,n}$,
d.\,h. die in dem $2m+3n$ kleiner als $\nu$ ist, schon berechnet
sind.
% [sic] Printed relative construction “die in dem ... ist”; grammatically expected “die, in denen ... ist”.

Man sieht daraus, daß die $a_{m,n}$ nur Potenzen von 3 im Nenner
enthalten. In den von H. A. \gesperrt{Schwarz} herausgegebenen
\glqq Formeln und Lehrsätzen zum Gebrauch der elliptischen
Funktionen\grqq, nach Vorlesungen von \gesperrt{Weierstrass}
(2. Ausgabe 1893), sind die $a_{m,n}$ bis zu $\nu=17$ berechnet. Es
zeigt sich dabei, daß diese Koeffizienten nicht nur ganze Zahlen
sind, sondern daß sie auch mit $\nu$ wachsende Potenzen von 3 als
Faktoren enthalten. Einen Beweis hierfür vermag ich nicht zu geben.

\begin{center}
\rule{2cm}{0.4pt}
\end{center}
